\section{Extremal families and the effect of separation}\label{sec:bundle-families}

We first state the stable path, end-attachment and three-arm families
needed for the extremum and the general subdivision theorem. Their
complete degree calculations and uniform estimates appear in
Appendices~\ref{sec:paths}--\ref{sec:spiders}. The distinction between
the two uses of these results is useful: paths and attachments attain
the numerical bound, whereas the quantitative three-arm estimate is
the initial case for arbitrarily many branching vertices.

\subsection{Paths}

Take the graph of Section~\ref{sec:graphs} to be a path with vertices
$0,\ldots,j$ and edges $i\colon i\to i+1$, $0\le i<j$. Put $b_c=2$ at
the interior vertices and $b_0,b_j\in\{1,2,3\}$ at the ends. Denote the
resulting variety by $X_j^{(b_0,b_j)}$. Its dimension and Picard number are
\begin{equation}\label{path:dimension}
 n=4j-2+b_0+b_j,\qquad \rho=5j+1.
\end{equation}
The Fano condition holds at every vertex. We write $X_j=X_j^{(1,1)}$.

\begin{theorem}\label{path:stable}
For every $j\ge1$ and every $b_0,b_j\in\{1,2,3\}$, the tangent bundle of
$X_j^{(b_0,b_j)}$ is slope stable with respect to its anticanonical polarization.
\end{theorem}


The proof is in Appendix~\ref{sec:paths}.

\subsection{Attainment in every dimension}

Put $r=(1,0)$ and $q=(0,1)$ in the hexagon lattice and define
\[
 A_0=(-r),\quad A_1=(q,-r-q),\quad A_2=(-r,q,-q),\quad
 A_3=(r,-r,q,-r-q).
\]
For $j\ge1$, take $j$ hexagon factors. Between hexagons $i-1$ and $i$
place a $\PP^2$ base factor with twisting vector $r_{i-1}-r_i$.
At the first hexagon attach one $\PP^1$ base factor for each root in
$A_a$; at the last attach one for each root in $-A_b$. Use the locally
trivial toric bundle fan over the product of these projective spaces,
with fibre $S_6^j$, and denote the resulting variety by $X_j(a,b)$.
\begin{theorem}\label{attach:stable}
For every $j\ge1$ and
\[
 (a,b)\in\{(1,0),(2,0),(1,1),(3,0),(1,2)\},
\]
$X_j(a,b)$ is smooth Fano and its tangent bundle is anticanonically stable.
Moreover,
\[
 \dim X_j(a,b)=4j+a+b,\qquad \rho(X_j(a,b))=5j+a+b+1.
\]
\end{theorem}


The proof is in Appendix~\ref{sec:attachments}.

\begin{corollary}[Sharpness in the bundle class]\label{attach:sharp}
For every $n\ge4$, the largest Picard number of a smooth toric Fano
$n$-fold with stable tangent bundle among the locally trivial toric
$S_6^j$-bundles over products of positive-dimensional projective spaces,
with the fan presentation~\eqref{interact:rays}, is
\[
 \left\lfloor\frac{5n}{4}\right\rfloor+1.
\]
\end{corollary}
\begin{proof}
Theorem~\ref{interact:bound} gives the upper bound. Write $n=4j+r$,
$0\le r\le3$. For $r=0$ use the ordinary path $X_j$ from
Theorem~\ref{path:stable}. For $r=1,2,3$ use respectively the pairs
$(1,0),(2,0),(3,0)$ in Theorem~\ref{attach:stable}. Their Picard
numbers are $5j+r+1=\lfloor5n/4\rfloor+1$.
\end{proof}

\subsection{The first branching case}

Stability at the Picard bound is compatible with branching, even when
the branches have different lengths. Let $\Gamma_{a,b,c}$ be the tree
with one vertex of degree three and arms containing $a,b,c$ edges.
Put $b_v=\deg(v)$ and let $Y_{a,b,c}=X(\Gamma_{a,b,c})$ be the graph
bundle of Section~\ref{sec:graphs}. Its base has one factor $\PP^3$,
three factors $\PP^1$, and $a+b+c-3$ factors $\PP^2$.

\begin{theorem}\label{spider:stable}
For all integers $a,b,c\ge2$, the tangent bundle of $Y_{a,b,c}$ is
stable with respect to $-K_{Y_{a,b,c}}$. Moreover,
\[
 \dim Y_{a,b,c}=4(a+b+c),\qquad
 \rho(Y_{a,b,c})=5(a+b+c)+1.
\]
\end{theorem}

Thus branching examples attain the bundle bound in every dimension
divisible by four and at least twenty-four. Their fans are not
lattice-isomorphic to the ordinary path fans in the same dimension:
the central $\PP^3$ rays form a primitive collection of size four,
whereas the path primitive collections have sizes two or three by
Lemma~\ref{graph:split-collections}.

The Fano property and connected ray matroid follow from
Proposition~\ref{graph:geometry}. The proof keeps independent subspace
choices on the three arms. A cubic integral gives their interaction at
the center. Contraction then bounds the effect of changing one arm
along the others, uniformly even when their lengths are very different.


The proof is in Appendix~\ref{sec:spiders}.

The contraction estimate behind the proof explains why separation is
effective. At a
degree-two vertex, integration over the base simplex acts on a positive
function by the operator
\[
 (\mathcal T f)(a')=
 \int_{-1}^1 f(a)\frac{(3+a-a')^2}{2}(2-|a|)\,da.
\]
Thus a path applies the same positive integral operator repeatedly.
For positive continuous functions define
\[
 d_H(f,g)=\log\left(\sup_a\frac{f(a)}{g(a)}
                            \sup_a\frac{g(a)}{f(a)}\right).
\]
Lemma~\ref{path:contraction} proves that for arbitrary positive
continuous initial functions and every $i\ge1$,
\begin{equation}\label{tree:boundary-contraction}
 d_H(\mathcal T^if,\mathcal T^ig)
 \le\log(81/25)(2/7)^{i-1}.
\end{equation}
The corresponding ray degrees are ratios of positive integrals. For
each hexagon or interior-base ray type, with left and right chain
indices $i,k\ge1$, its normalized degree $\delta(i,k)$ consequently
satisfies
\[
 |\log\delta(i,k)-\log\delta_\infty|
 \le\log(81/25)\bigl((2/7)^{i-1}+(2/7)^{k-1}\bigr).
\]
The limiting degree is independent of the distant end data. An end-ray
degree has only a one-sided limit and can retain dependence on its own
end factor. These are convergence estimates; they do not assert
monotonicity in the path lengths.


Stability also requires a uniform bound for the deficits of all
relevant subspaces. Section~\ref{sec:subdivided} supplies the rank
formula and the gluing argument that preserve those bounds when smaller
trees are joined.
